\section{Kim Polese}
Liderou o lançamento do Java na Sun e depois cofundou a Marimba (fonte:
\citeobra{bio:polese-chm}).

\begin{itemize}
\obra*{artigo:polese-talks-java}{Entrevista de 1996, no último dia dela na
  JavaSoft: por que a Sun \enquote{deu de graça} o Java, o licenciamento e
  os mercados que ela via para a linguagem.}
\item No documentário \citetitle{video:java-story}, como a primeira gerente
  de produto do Java.
\end{itemize}
Não achamos outro texto ou palestra dela sobre o Java que pudéssemos
conferir.

\section{James Duncan Davidson}
Na Sun, foi responsável pelas especificações Servlet~2.1/2.2 e JAXP~1.0; é
o autor original do Tomcat e do Ant (fonte: \citeobra{bio:davidson-ant}).

\begin{itemize}
\obra{doc:ant-faq}{O FAQ do Ant (texto do projeto, não dele) conta que o
  Ant nasceu dentro do Tomcat, para compilá-lo, e que o nome vem de
  \enquote{Another Neat Tool}, segundo Davidson. Veio antes do Maven e do
  Gradle.}
\item No documentário \citetitle{video:java-story}, sobre o Tomcat, o Ant e
  o Java no servidor.
\end{itemize}
Material técnico escrito por ele mesmo, aberto e conferido, não achamos.

\section{Carla Schroer}
Foi Director of Java Compatibility na Sun, segundo o Foojay (fonte:
\citeobra{artigo:foojay-java-story}). Não achamos página oficial dela.

\begin{itemize}
\item No documentário \citetitle{video:java-story}, sobre compatibilidade
  (o \enquote{write once, run anywhere}) e a disputa com a Microsoft.
\end{itemize}
Nenhuma palestra ou texto técnico dela foi encontrado e confirmado; o
documentário é, por enquanto, a fonte.
